%=============================================================================
% MODULE 04: LITERATURE REVIEW (Additional Expansion)
%=============================================================================
% Frontiers in Public Health — Thematic narrative-critical review
%   Block 1: Algorithmic bias in healthcare
%   Block 2: Diabetes disparities
%   Block 3: Fairness metrics
%   Block 4: Mitigation strategies
%   Block 5: Intersectionality
%=============================================================================

\section{Literature Review}

\subsection{Algorithmic Bias in Healthcare}

Algorithmic bias has been documented across numerous healthcare applications, 
from risk prediction to diagnostic imaging \cite{rajkomar2018}. A landmark 
study by Obermeyer et al. \cite{obermeyer2019} revealed that a widely used 
healthcare algorithm systematically underestimated the health needs of Black 
patients by using healthcare costs as a proxy for health status, failing to 
account for structural barriers to accessing care. This finding underscored 
how seemingly objective algorithmic decisions can embed and scale historical 
inequities, affecting millions of patients annually.

The Obermeyer study demonstrated that at a given risk score, Black patients 
were considerably sicker than White patients, as evidenced by conditions 
such as diabetes, hypertension, and chronic kidney disease. The algorithm 
used healthcare costs as a proxy for health needs, but because Black patients 
historically had less access to healthcare and thus lower costs, the model 
systematically underestimated their health needs. This case illustrates how 
historical inequities can be perpetuated and amplified through algorithmic 
decision-making.

In diabetes specifically, Chen et al. \cite{chen2019} demonstrated that 
prediction models trained on electronic health records showed differential 
performance across racial groups, with higher false-negative rates for 
minority populations. The study found that models trained predominantly on 
data from White patients showed reduced accuracy when applied to Black and 
Hispanic patients, suggesting that training data composition significantly 
influences model fairness.

Similarly, a systematic review by Rajkomar et al. \cite{rajkomar2018} 
identified performance disparities in clinical ML models across demographic 
subgroups, particularly in settings with imbalanced representation. The 
review highlighted that many clinical ML studies fail to report performance 
metrics stratified by demographic groups, making it difficult to assess 
fairness in existing literature.

The sources of algorithmic bias are multifaceted and interconnected 
\cite{mehrabi2021}:

\begin{itemize}[leftmargin=*]
    \item \textbf{Historical bias}: Past discrimination embedded in training 
          data, reflecting unequal healthcare access and treatment patterns. 
          This includes differential treatment protocols, varying levels of 
          care, and systemic barriers that have historically marginalized 
          certain populations.
    \item \textbf{Representation bias}: Underrepresentation of certain groups 
          in training datasets, leading to poor generalization for these 
          populations. When models are trained on data that does not adequately 
          represent the diversity of the target population, they may fail to 
          capture the complex relationships that exist in underrepresented 
          groups.
    \item \textbf{Measurement bias}: Inconsistent data collection protocols 
          across populations, introducing systematic errors. This can include 
          differences in measurement tools, protocols, and standards that 
          vary across clinical settings and populations.
    \item \textbf{Aggregation bias}: One-size-fits-all models that ignore 
          meaningful population differences in disease presentation and 
          progression. Different populations may exhibit distinct disease 
          phenotypes, comorbidity patterns, and treatment responses that 
          are not captured by aggregated models.
\end{itemize}

The consequences of algorithmic bias extend beyond individual patient harm 
to affect public health at scale. When biased algorithms are deployed in 
clinical decision support systems, they can systematically disadvantage 
entire populations, exacerbate existing health disparities, and undermine 
trust in healthcare systems. The cumulative effect of these biases can 
create feedback loops where disadvantaged populations receive worse care, 
leading to poorer outcomes that further reinforce the biases in the data.

\subsection{Diabetes Disparities}

Diabetes disproportionately affects marginalized communities due to social 
determinants of health \cite{narayan2017}. The International Diabetes 
Federation \cite{idf2021} reports that low- and middle-income countries bear 
the greatest burden, with prevalence rates exceeding those in high-income 
settings by 2--3 fold. Within high-income countries, racial and ethnic 
minorities consistently experience higher prevalence and worse outcomes.

The Centers for Disease Control and Prevention \cite{cdc2020} reports 
significant racial and ethnic disparities in the United States, as summarized 
in Table~\ref{tab:diabetes_disparities}. These disparities reflect a complex 
interplay of social, economic, and environmental factors that create unequal 
health outcomes across populations.

\begin{table}[H]
\centering
\caption{Diabetes Prevalence by Race/Ethnicity in the United States 
(Source: CDC, 2020)}
\begin{tabular}{lcc}
\toprule
\textbf{Group} & \textbf{Prevalence (\%)} & \textbf{Relative Risk} \\
\midrule
White & 7.4 & 1.00 \\
Black & 11.7 & 1.58 \\
Hispanic & 12.5 & 1.69 \\
Asian & 8.0 & 1.08 \\
American Indian & 14.7 & 1.99 \\
\bottomrule
\end{tabular}
\label{tab:diabetes_disparities}
\end{table}

These disparities are driven by multiple interconnected factors 
\cite{hill2022}: unequal access to healthcare services, higher rates of 
obesity and physical inactivity, food insecurity and limited access to 
healthy foods, chronic stress from discrimination, and lower socioeconomic 
status. When ML models are trained on data reflecting these disparities 
without appropriate correction, they may learn and perpetuate these inequities, 
creating a feedback loop that further marginalizes already disadvantaged 
populations.

The social determinants of health that drive diabetes disparities are deeply 
rooted in structural inequities, including residential segregation, 
educational disparities, employment discrimination, and historical trauma. 
These factors create cumulative disadvantage that manifests in higher 
diabetes prevalence, later diagnosis, and poorer outcomes among marginalized 
populations. Algorithmic models trained on data shaped by these structural 
factors risk encoding and amplifying these inequities.

\subsection{Fairness Metrics}

Multiple fairness metrics have been proposed to quantify algorithmic bias 
\cite{corbett2018}. The most commonly used metrics include:

\subsubsection{Equalized Odds}

Equalized Odds, introduced by Hardt et al. \cite{hardt2016}, requires that 
the true positive rate (sensitivity) and false positive rate (1 - specificity) 
are equal across groups:

\begin{equation}
P(\hat{Y}=1 \mid Y=1, A=a) = P(\hat{Y}=1 \mid Y=1, A=b) \quad \forall a,b
\label{eq:eo_tpr}
\end{equation}

\begin{equation}
P(\hat{Y}=1 \mid Y=0, A=a) = P(\hat{Y}=1 \mid Y=0, A=b) \quad \forall a,b
\label{eq:eo_fpr}
\end{equation}

\noindent where $\hat{Y}$ is the predicted outcome, $Y$ is the true outcome, 
and $A$ is the sensitive attribute. Equalized Odds is particularly relevant 
in clinical settings where both missed diagnoses (false negatives) and 
unnecessary interventions (false positives) carry significant consequences.

The Equalized Odds criterion is especially important in diabetes screening 
because both types of errors have serious clinical implications. False 
negatives can lead to delayed treatment and increased risk of complications, 
while false positives can result in unnecessary anxiety, additional testing, 
and potential overtreatment. By requiring equal error rates across groups, 
Equalized Odds ensures that no population bears a disproportionate burden 
of these errors.

\subsubsection{Demographic Parity}

Demographic Parity requires that the prediction is statistically independent 
of the sensitive attribute \cite{dwork2012}:

\begin{equation}
P(\hat{Y}=1 \mid A=a) = P(\hat{Y}=1 \mid A=b) \quad \forall a,b
\label{eq:dp}
\end{equation}

While simpler to implement, Demographic Parity does not account for 
legitimate differences in base rates across groups, which may lead to 
paradoxical conclusions. For example, if one group has a genuinely higher 
prevalence of diabetes, requiring equal prediction rates could lead to 
under-diagnosis in that group.

The tension between Demographic Parity and Equalized Odds highlights the 
complexity of defining fairness in clinical contexts. Different fairness 
criteria may be appropriate depending on the specific clinical application, 
the consequences of different types of errors, and the underlying 
prevalence patterns across populations.

\subsubsection{Calibration}

A model is well-calibrated if the predicted probabilities match the true 
outcome rates across groups \cite{chouldechova2017}:

\begin{equation}
P(Y=1 \mid \hat{P}=p, A=a) = P(Y=1 \mid \hat{P}=p, A=b) \quad \forall a,b,p
\label{eq:cal}
\end{equation}

Chouldechova \cite{chouldechova2017} demonstrated that except in trivial 
cases, it is impossible to simultaneously satisfy calibration, equalized 
odds, and predictive parity, necessitating careful consideration of which 
fairness criterion is most appropriate for the clinical context.

Calibration is particularly important in clinical decision-making because 
it ensures that predicted probabilities are meaningful and actionable. 
A well-calibrated model allows clinicians to interpret risk scores 
consistently across patient groups, enabling more equitable clinical 
decision-making. However, achieving calibration across groups may require 
group-specific adjustments that raise additional fairness considerations.

\subsection{Bias Mitigation Strategies}

Bias mitigation strategies can be categorized by their position in the 
ML pipeline \cite{friedler2019}:

\begin{enumerate}[leftmargin=*]
    \item \textbf{Pre-processing} (data-level): Modify training data to 
          remove bias before model training. Techniques include reweighting 
          (assigning different weights to samples), resampling (oversampling 
          underrepresented groups), and data augmentation (generating synthetic 
          samples using SMOTE or similar methods). Pre-processing methods are 
          attractive because they are model-agnostic and can be applied to any 
          dataset without modifying the learning algorithm.
    
    \item \textbf{In-processing} (algorithm-level): Modify the learning 
          algorithm to incorporate fairness constraints. Approaches include 
          adversarial debiasing (removing sensitive information from learned 
          representations), fairness-constrained optimization, and 
          regularization methods that penalize unfair predictions. In-processing 
          methods offer the advantage of jointly optimizing for performance 
          and fairness, but require access to the learning algorithm.
    
    \item \textbf{Post-processing} (prediction-level): Modify model 
          predictions to achieve fairness. Methods include threshold adjustment 
          (applying different decision thresholds per group), calibration 
          (equalizing predicted probabilities), and reject option (abstaining 
          from prediction on uncertain cases). Post-processing methods are 
          flexible and can be applied to any trained model, but may not address 
          the root causes of bias.
\end{enumerate}

Each category offers distinct advantages and trade-offs. Pre-processing 
methods are model-agnostic but may reduce overall accuracy. In-processing 
methods optimize fairness and accuracy jointly but require algorithmic 
modification. Post-processing methods are flexible but may not address 
the root causes of bias. In practice, a combination of approaches may be 
most effective, as different sources of bias may require different 
mitigation strategies.

\subsection{Intersectionality in Algorithmic Fairness}

Intersectionality theory, originally proposed by Crenshaw \cite{crenshaw1989}, 
highlights how multiple forms of discrimination interact to create unique 
modes of disadvantage. In algorithmic fairness, intersectional analysis 
examines bias at the intersection of multiple sensitive attributes (e.g., 
race \textit{and} gender) rather than each attribute independently 
\cite{buolamwini2018}.

Buolamwini and Gebru \cite{buolamwini2018} demonstrated this phenomenon 
in facial recognition systems, where commercial classifiers exhibited 
significant performance disparities for darker-skinned women, despite 
reasonable accuracy for lighter-skinned individuals and men. This finding 
highlighted the limitations of single-attribute fairness analysis and the 
importance of examining intersectional effects.

In healthcare contexts, intersectional bias is particularly concerning 
because health disparities often cluster along multiple dimensions. 
Individuals who are simultaneously members of multiple marginalized groups 
may experience compounded disadvantage that is not captured by analyzing 
each attribute in isolation. For example, Black women may face both racial 
and gender discrimination in healthcare, creating unique barriers that 
are not captured by analyzing race or gender separately.

Despite its importance, intersectional analysis remains underrepresented 
in healthcare AI fairness research, representing a critical gap that this 
study addresses. The few existing studies that have examined intersectional 
effects in healthcare AI have found significant disparities at the 
intersection of race and gender, suggesting that single-attribute analyses 
may systematically underestimate the extent of algorithmic bias.

The intersectionality framework also highlights the importance of 
considering structural and contextual factors that shape health disparities. 
Algorithmic bias does not occur in isolation but is embedded in broader 
systems of inequality that influence data collection, model development, 
and deployment. Understanding these structural factors is essential for 
developing effective strategies to address algorithmic bias in healthcare.
